\section{Introduction}\label{sec:introduction}

Bitcoin \cite{bitcoin} is usually described as a ledger of coins. Ordinal theory gives that
ledger a second reading: every satoshi has a number, sats move through
transactions in first-in-first-out order, and a particular sat can therefore
be followed across the chain. HOPPER uses that identity directly. A frog is
not a token contract, a database row, or an entry in a server's table. A frog
is one satoshi, and its public life is the sequence of transactions that
carry that satoshi.

The game is small enough to state in one paragraph. A player mints a frog by
inscribing a self-contained art document on a sat. A one-time \textsc{prepare}
may move the frog from its mint output into the custody output. To move the
frog forward, the owner spends the sat through one of three Taproot leaves and
publishes a 55-byte state inscription naming the direction it picked: left,
straight or right. The crossing confirms at height $T$. At block $T+1$ the
block hash wraps to exactly one direction for the whole block, and every
frog whose crossing confirmed in $T$ that picked that direction dies at
$T+\kappa$. A frog that survives enough crossings reaches \emph{home} and
becomes immortal. A plain transfer moves the sat without publishing anything
and without changing the frog's state. The chain is the referee, and the
owner chooses \emph{when} to hop and \emph{which way}, never whether the road
agrees.

\subsection{Design goals}

The construction is organized around four goals.

\begin{description}[leftmargin=1.4em,style=nextline,itemsep=0.35em]
  \item[State on the sat, not beside it.]
  Every state transition is an ord inscription on the frog's own sat. A stock
  \code{ord} viewer can reconstruct the frog from chain data alone, without a
  HOPPER server. The indexer accelerates this; it is not the source of truth.

  \item[Custody by Bitcoin Script.]
  A frog in custody sits in a Taproot output whose internal key is a
  nothing-up-my-sleeve point. There is no key-path spend. Three leaves sign a
  crossing, one per direction, each carrying the state envelope; the fourth
  is a plain owner signature for transfers. A valid crossing is therefore
  both a signature and a protocol message, and the direction is bound into
  the leaf that authorized it.

  \item[Outcomes the author cannot choose.]
  At block $T+1$ the block hash selects one direction for the entire block.
  A crossing dies exactly when the direction committed in its body equals
  that wrap. The block hash does not exist when the crossing is signed, and
  the function is public. No author, owner or indexer can declare an outcome.
  A block producer can influence the candidate hash or withhold a block;
  that is miner influence, not player agency, and the protocol does not
  pretend otherwise.

  \item[An implementation small enough to audit.]
  The consensus surface is a 55-byte body, a 34-byte send leaf, a 161-byte
  crossing leaf, a 36-byte parameter encoding, and three integer functions:
  generation, required crossings, and the block wrap. The reference
  implementation is split into a pure protocol crate, a pure state machine,
  and a network-facing API.
\end{description}

\subsection{Why a satoshi?}

A satoshi is the smallest unit Bitcoin can own, and ordinal theory gives it a
stable name. This has three consequences that matter for a game.

First, ownership is UTXO ownership. There is no separate balance to reconcile
and no issuer-controlled registry that can reassign a frog. Second, the sat
carries its history. A state inscription is not attached to a mutable object;
it is attached to the same sat as the mint, so the child relationship is
cryptographically anchored by the satpoint as well as by the parent tag.
Third, the identity survives transfers. A frog can be sent to an ordinary
address, held, and later brought back into custody without losing its sat or
its accumulated history.

The price of this design is that a frog is only as movable as the UTXO that
contains it. A wallet must know which ordinal ranges a coin carries before it
spends that coin. The protocol makes this requirement explicit: a coin whose
ranges cannot be attributed is still counted in the balance, but it is shown
as unspendable rather than spent blindly.

\subsection{What is new in v8}

The game rules introduced by wire version \code{0x07} still shape
everything that follows. Wire version \code{0x08} adds a proof-bound
envelope, a proof-checking fold and a single legacy exception.

\begin{enumerate}[leftmargin=1.6em,itemsep=0.3em]
  \item \textbf{The road asks a question, not a die roll.} Each crossing body
        carries a direction: left, straight or right. The next block's hash
        wraps to one direction and applies it to every crossing in the
        previous block. There is no per-Hopper death table and no per-Hopper
        randomness.
  \item \textbf{Custody has four leaves.} Three crossing leaves, one per
        direction, plus one send leaf. Choosing a direction selects a leaf
        whose bytes were fixed when the state output was created.
  \item \textbf{Difficulty runs in generations.} A generation starts at five
        required crossings, adds one every 144 blocks up to twenty-five,
        holds for 3{,}600 blocks, and resets. Each generation stamps its frogs
        \code{G1}, \code{G2}, and so on; the label is fixed for life, the
        target is not.
\end{enumerate}

Version 8 changes the bytes and the verification story, not the road. The
mint body is now \code{live_v8_html}, the state magic is
\code{FRG8CROSS}/\code{0x08}, and the embedded fold checks every
body-valid child against the raw transaction that published it. A proof that
is missing or false makes the fold report \emph{unverified}, rendered as
\code{NAME} \roadglyph{} \code{---}; it never prints a count or claims
\code{HOME} on incomplete evidence. The fold ignores a body-invalid child,
but a body-valid child with a failed proof is a hard stop.

Exactly one released v7 Hopper is carried across the cut. Its sat is the
compiled \code{GRANDFATHERED_SAT}. For that sat only, the v7 mint envelope
and digest, the \code{V7} body and the v7 \code{FRG7CROSS}/\code{0x07}
state path are accepted, because already-confirmed v7 custody cannot be
rewritten. Every other sat must use v8. The released v7 card's fold predates
the proof check; it is grandfathered for custody and play, not as an
authoritative art oracle. A proof-checked view comes from the v8 app or
indexer.

The mint reveals a self-contained HTML document whose canonical fold reads
only relative \code{/r/} and \code{/content/} endpoints on the viewer's own
ord node. A crossing publishes an \code{application/octet-stream} body
containing the frog's sat, mint identity and chosen direction, together with
an ord parent tag. The first state output must be created by a later
transaction --- the first transfer or a one-time \textsc{prepare} --- because
the mint transaction cannot embed its own txid; after that, each crossing
spends the previous state output and creates the next one.

\subsection{Contributions and outline}

This paper makes the following contributions.

\begin{enumerate}[leftmargin=1.6em,itemsep=0.25em]
  \item A precise inscription format for mint and crossing, including the ord
        envelope, the parent encoding, the canonical PNG encoder, the single
        accepted v8 mint-body digest and its one legacy exception
        (Section~\ref{sec:inscriptions}).
  \item A consensus state machine for mint, hop, settlement, transfer and
        burn, with a parameter commitment, a generation schedule and a
        block-wide wrap (Sections~\ref{sec:parameters} and~\ref{sec:game}).
  \item A four-leaf Taproot construction that binds the protocol message and
        the chosen direction to the spend, with exact leaf scripts, control
        blocks, witnesses and transaction shapes
        (Section~\ref{sec:custody}).
  \item A standalone Rust indexer with deterministic state, snapshots, a
        reorg window and a documented HTTP API
        (Section~\ref{sec:indexer}).
  \item A JavaScript wallet/PWA that keeps keys in a WebCrypto vault,
        reproduces the interface and sound, and exports emergency descriptors
        for its staged outputs (Section~\ref{sec:wallet}).
  \item A proof-checking fold, a fail-closed unverified state and the single
        grandfathered-v7 rule, plus a security analysis and a
        conformance-oriented test plan
        (Sections~\ref{sec:inscriptions}, \ref{sec:security} and
        \ref{sec:validation}).
\end{enumerate}

Section~\ref{sec:model} fixes notation and the trust model.
Section~\ref{sec:parameters} states the constants, the generation schedule
and the wrap. Section~\ref{sec:inscriptions} specifies the envelopes.
Section~\ref{sec:game} gives the state machine.
Section~\ref{sec:custody} gives the tapscript tree and transaction shapes.
Sections~\ref{sec:indexer} and~\ref{sec:wallet} describe the reference
infrastructure. Section~\ref{sec:security} analyses attacks and limits, and
Section~\ref{sec:validation} lists the test obligations. The appendices
reproduce the generation schedule and the deployment constants.

\begin{invariant}{The one-sentence rule}
A Hopper is a satoshi. Its state is the ordered set of valid HOPPER
inscriptions on that satoshi. Its fate is a function of the chain, not of
its owner.
\end{invariant}
